% ============================================================
% Civilization-ABM v3 — Manuscript
% Population Scale, Interaction Memory, and Multidimensional
% Inequality in an Agent-Based Civilization Model
% ============================================================
\documentclass[12pt, letterpaper]{article}

\usepackage[margin=1in]{geometry}
\usepackage{amsmath, amssymb}
\usepackage{booktabs}
\usepackage{graphicx}
\usepackage{hyperref}
\usepackage{natbib}
\usepackage{setspace}
\usepackage{parskip}
\usepackage{array}
\usepackage{multirow}
\usepackage{caption}

\doublespacing

\title{\textbf{Population Scale, Interaction Memory, and\\
Multidimensional Inequality in an Agent-Based\\
Civilization Model: A Replication and Extension Study}}

\author{
  Juan Moisés de la Serna Tuya\\
  Universidad Internacional de La Rioja (UNIR)\\
  Logroño, Spain\\
  \texttt{juanmoises.delaserna@unir.net}\\
  ORCID: 0000-0002-8401-8018
}

\date{}

\begin{document}

\maketitle

% ============================================================
\begin{abstract}
% ============================================================

This paper presents Civilization-ABM v3, a third-generation agent-based
model extending a validated framework with two novel mechanisms: (1)
\emph{interaction memory}, a circular-buffer system inspired by large-scale
swarm intelligence architectures \citep{mirofish2025, yang2024oasis},
through which agents record the outcomes of their last five interactions
and weight strategic imitation by the reputation of the role model
(indirect reciprocity; \citealp{nowak1998indirect}); and (2) a
\emph{multidimensional inequality measurement suite} comprising the Gini
coefficient, Theil T index, Palma ratio, top-1\% wealth share, top-10\%
share, and social mobility rate. The model scales from N~=~500 agents on
a 35$\times$35 landscape (v2) to N~=~1,000 agents on a 50$\times$50
landscape, providing the first population-scale robustness test of the
Civilization-ABM findings. Across 48 experimental conditions and 1,440
replications (30 per condition; 1,500 steps each), eight research
questions are addressed. Key findings include: (1) Population scaling
from N\,=\,500 to N\,=\,1,000 reduces equilibrium Gini by 44.9\%
(0.361~$\to$~0.199), but fiscal policy effectiveness collapses due to
endemic institutional failure (100\% collapse rate); (2) competitive
strategy dominance is universal across 38 of 48 conditions, replicating
the central v2 finding; (3) institutional collapse is pervasive,
occurring in 94.2\% of simulation steps under baseline conditions;
(4) wealth floor mechanisms produce the largest inequality reductions
(progressive + floor: Gini\,=\,0.063, $-$68\% vs.\ baseline; welfare
state: Gini\,=\,0.024), with stable regimes in 53--87\% of replications;
(5) interaction memory has negligible marginal effect on inequality
(all memory--no-memory Gini differences $\leq$\,0.007) but modestly
increases cooperative share in combined floor + memory conditions;
(6) World Bank calibration produces qualitatively correct ordering
(wb\_nordic: 0.001; wb\_south\_africa: 0.638) but overshoots targets for
most archetypes due to the model's binary institutional collapse.
All source code is released under MIT license at
\texttt{https://github.com/juanmoisesd/civilization-abm-v3}.

\medskip
\noindent\textbf{Keywords:} agent-based modeling; wealth inequality;
interaction memory; replicator dynamics; Theil index; Palma ratio;
social mobility; swarm intelligence; Sugarscape; institutional collapse

\end{abstract}

% ============================================================
\section{Introduction}
% ============================================================

The emergence and persistence of wealth inequality constitutes one of
the central unresolved puzzles of social science \citep{atkinson2015,
piketty2014}. Agent-based modeling (ABM) offers a uniquely suited
methodological approach: by simulating populations of heterogeneous,
interacting agents, ABM allows researchers to observe how complex
distributional patterns emerge from simple local rules --- the
\emph{generative} approach to social science \citep{epstein1996,
axelrod1997}.

Two companion papers established the Civilization-ABM framework.
\citet{delaserna2024v1} introduced Civilization-ABM v1, demonstrating
that progressive fiscal policy reduces the Gini coefficient by 13.3\%
relative to a no-taxation baseline, while high initial wealth dispersion
paradoxically reduces equilibrium inequality through a redistribution
amplification mechanism. \citet{delaserna2025v2} introduced
Civilization-ABM v2, extending the framework with evolutionary strategy
learning via replicator dynamics, spatial resource heterogeneity,
institutional collapse and recovery, and empirical calibration against
World Bank Gini data. Across 41 conditions and 2,460 replications
(N~=~500 agents), five findings emerged: competitive strategy dominance
is universal and invariant to fiscal redistribution; geographic
concentration exerts minimal effect on inequality; institutional
collapse is endemic; a floor trap mechanism produces extreme inequality;
and World Bank calibration fails systematically due to the model's 100\%
redistribution efficiency assumption.

These findings, however, were obtained at a population scale of
N~=~500 agents on a 35$\times$35 landscape --- a scale chosen for
computational tractability with Mesa 3.x \citep{masad2015mesa}, not for
behavioral realism. Empirically relevant societies contain orders of
magnitude more participants, and network effects in social learning scale
non-linearly with population size \citep{barabasi1999}. Whether the
findings of v2 constitute genuine social regularities or artifacts of
small-scale simulation remains untested.

Simultaneously, a new generation of large-scale social simulation
architectures has emerged. MiroFish \citep{mirofish2025}, built on the
OASIS engine \citep{yang2024oasis}, demonstrated that swarm simulations
with thousands to one million agents can produce qualitatively distinct
emergent dynamics invisible at smaller scales. A key architectural
innovation in these systems is \emph{agent memory}: each agent maintains
a record of past interactions that modulates future behavioral choices.
AgentTorch \citep{chopra2024agenttorch} further demonstrated that
population-scale ABMs (up to 8.4 million agents) can be calibrated
against real-world data, establishing a new methodological standard for
computational social science.

These developments motivate three extensions in Civilization-ABM v3:
(1) \emph{population scaling}, testing whether v2 findings replicate at
N~=~1,000 agents on a 50$\times$50 landscape; (2) \emph{interaction
memory}, implementing a MiroFish-inspired circular buffer of five
interaction outcomes per agent, with reputation-weighted replicator
dynamics (indirect reciprocity; \citealp{nowak1998indirect,
santos2008social}); and (3) \emph{multidimensional inequality
measurement}, replacing the Gini-only measurement of v2 with a full
suite including Theil T, Palma ratio, top-1\% share, and social
mobility.

Eight research questions guide the analysis:

\begin{itemize}
  \item \textbf{RQ1:} Do the five principal findings of Civilization-ABM
    v2 replicate when population size doubles from N~=~500 to
    N~=~1,000?
  \item \textbf{RQ2:} Does population scale alter the competitive
    strategy equilibrium identified in v2?
  \item \textbf{RQ3:} Does the 50$\times$50 landscape amplify or
    attenuate spatial inequality effects relative to the 35$\times$35
    landscape?
  \item \textbf{RQ4:} Does endemic institutional collapse persist at
    population scale, and how does its frequency compare to v2
    baseline?
  \item \textbf{RQ5:} Do the Theil T index and Palma ratio reveal
    distributional patterns invisible to the Gini coefficient?
  \item \textbf{RQ6:} Does interaction memory --- via reputation-weighted
    replicator dynamics --- alter the competitive equilibrium identified
    in v2?
  \item \textbf{RQ7:} Does memory reduce the floor trap effect observed
    when wealth floors are implemented without redistribution?
  \item \textbf{RQ8:} Is social mobility reduced during institutional
    collapse, and does memory-based reputation increase it?
\end{itemize}


% ============================================================
\section{Background and Related Work}
% ============================================================

\subsection{Population Scale in Agent-Based Social Simulation}

The relationship between population size and emergent inequality
dynamics in ABM is theoretically non-trivial. Network effects in
social learning --- the propagation of behavioral strategies through
social contacts --- scale non-linearly with N in both small-world and
scale-free topologies \citep{watts1998, barabasi1999}. In the original
Sugarscape model \citep{epstein1996}, population size was found to
influence the shape of wealth distributions, though the qualitative
Pareto structure persisted across scales. More recently,
\citet{chopra2024agenttorch} demonstrated that population-scale ABMs
(N $\sim$ 10$^6$ -- 10$^7$) produce emergent coordination phenomena
absent at N $\sim$ 10$^3$, motivating the explicit study of
scale-dependence in validated ABM frameworks.

The present paper provides a direct N-scaling test: all mechanisms
from Civilization-ABM v2 are preserved unchanged, and only population
size (500 $\to$ 1,000) and landscape dimensions
(35$\times$35 $\to$ 50$\times$50) are varied.

\subsection{Interaction Memory and Indirect Reciprocity}

The role of memory in sustaining cooperative equilibria is one of the
most robust findings in evolutionary game theory. \citet{axelrod1984}
demonstrated that ``Tit-for-Tat'' --- a memory-one strategy --- can
outcompete defectors in repeated Prisoner's Dilemma tournaments.
\citet{nowak1998indirect} formalized \emph{indirect reciprocity}: agents
do not need direct experience with a partner; they can rely on the
partner's \emph{reputation} --- a social score derived from observed
interactions with third parties. This mechanism substantially expands
the conditions under which cooperation can be evolutionarily stable.

\citet{santos2008social} showed that social diversity in interaction
patterns --- heterogeneous memory depths and network positions ---
promotes cooperative emergence even when direct reciprocity would fail.
In the present model, memory is implemented as a circular buffer of five
interaction outcomes $o_t \in \{-1, 0, +1\}$, and reputation is updated
as:
\begin{equation}
  r_i(t+1) = \text{clip}\!\left(r_i(t) + 0.01 \cdot
    \bar{o}_i(t),\; 0, 2\right)
\end{equation}
where $\bar{o}_i(t)$ is the mean outcome over the five most recent
interactions. This creates a continuous reputation signal from the
history of cooperative, neutral, or competitive behavior.

The reputation-weighted replicator dynamic modifies the standard
probability of strategy adoption \citep{schuster1983}:
\begin{equation}
  p_{\text{copy}}^* = p_{\text{copy}} \cdot
    \underbrace{\left(0.5 + 0.5 \cdot \frac{r_{b^*}}{2}\right)}_{
    \text{reputation factor} \in [0.25, 1.0]}
\end{equation}
where $r_{b^*} \in [0, 2]$ is the reputation of the role model $b^*$.
Agents with high reputation are preferentially imitated regardless of
their momentary wealth advantage.

\subsection{Beyond Gini: Multidimensional Inequality Measurement}

The Gini coefficient, while widely used, is known to have differential
sensitivity across the wealth distribution: it is most sensitive to
transfers near the median and least sensitive to changes in the tails
\citep{atkinson1970}. A growing literature advocates for a suite of
complementary inequality measures \citep{palma2011}.

The \textbf{Theil T index} \citep{theil1967} is given by:
\begin{equation}
  T = \frac{1}{N} \sum_{i=1}^{N} \frac{w_i}{\bar{w}}
    \ln\!\left(\frac{w_i}{\bar{w}}\right)
\end{equation}
Theil T is particularly sensitive to inequality in the upper tail
(extreme wealth concentration), making it complementary to the Gini.

The \textbf{Palma ratio} \citep{palma2011} is defined as the ratio of
the wealth share of the richest 10\% to the wealth share of the poorest
40\%:
\begin{equation}
  P = \frac{W_{90\text{--}100}}{W_{0\text{--}40}}
\end{equation}
\citet{palma2011} documented that the middle 50\% of income
distributions are remarkably stable across countries, while the
distributional contest occurs almost entirely between the top decile and
the bottom four deciles --- making the Palma ratio an informative
measure of structural polarization.

\textbf{Social mobility} is operationalized as the fraction of alive
agents transitioning between social class assignments (lower / middle /
upper) per simulation step. This provides a dynamic measure of
distributional fluidity complementary to the static inequality indices
\citep{chetty2014}.

\subsection{Swarm Intelligence Architectures and MiroFish}

Recent advances in large-scale multi-agent AI simulation have produced
architectures with qualitatively different properties from classical
ABM. MiroFish \citep{mirofish2025}, built on the OASIS engine
\citep{yang2024oasis}, implements swarm simulations with up to one
million agents, each maintaining persistent memory and behavioral
history. The key insight from MiroFish is that collective prediction
accuracy improves non-linearly with agent count and memory depth,
suggesting that memory-enhanced agents at population scale may produce
emergent phenomena absent in memoryless small-scale models. While
MiroFish targets opinion dynamics and social media simulation rather
than economic inequality, its memory architecture motivates the
analogous mechanism implemented here.


% ============================================================
\section{Model and Methods}
% ============================================================

\subsection{Model Architecture}

Civilization-ABM v3 is implemented in Python 3.12 using NumPy-vectorized
operations (eliminating Mesa's per-agent iteration overhead, yielding a
10$\times$ speedup over v2) and multiprocessing parallelism via Python's
\texttt{multiprocessing.Pool}. The model comprises six coupled modules:
(1) agents, (2) resource landscape, (3) evolutionary dynamics,
(4) social network, (5) institutional system, and (6) data collection.
All mechanisms from v2 are preserved unchanged; v3 introduces two
extensions: interaction memory (Section~\ref{sec:memory}) and
multidimensional inequality measurement. Full source code is available
under MIT license at \texttt{https://github.com/jmdelaserna/
civilization-abm-v3}.

\subsection{Agents}

Each simulation initializes N~=~1,000 agents with the following
attributes:
\begin{itemize}
  \item \textbf{Wealth} ($w_i$): drawn from LN(2.3, $\sigma_0$) where
    $\sigma_0$ is the initial inequality parameter
  \item \textbf{Strategy} ($s_i$): drawn uniformly from
    \{cooperative, competitive, neutral\}
  \item \textbf{Metabolism} ($m_i$): drawn from
    LN(0, 0.6)$\times m_{\text{base}}$, generating heterogeneous
    structural costs
  \item \textbf{Vision} ($v_i$): drawn from \{1, 2, 3\} with weights
    \{0.50, 0.35, 0.15\}
  \item \textbf{Position} $(x_i, y_i)$: placed randomly on unoccupied
    landscape cells
  \item \textbf{Reputation} ($r_i$): initialized at 1.0
  \item \textbf{Memory} ($\mathbf{o}_i$): circular buffer of 5
    interaction outcomes, initialized at zero (new in v3)
\end{itemize}

\subsection{Resource Landscape}

The landscape is a 50$\times$50 discrete grid (expanded from 35$\times$35
in v2 to maintain agent density $\approx$ 0.4 agents/cell). Resource
capacity at each cell follows:
\begin{equation}
  C(x,y) = \text{clip}\!\left(1 + \sum_{k=1}^{K} s_k \exp\!\left[
    -\frac{(x-p_{k,x})^2 + (y-p_{k,y})^2}{2\sigma_k^2}
  \right], 1, C_{\max}\right)
\end{equation}
with $K$ Gaussian peaks, $C_{\max} = 20$, default $K = 3$, $\rho = 0.5$.
Agents move toward the free adjacent cell with the highest current
resources within their vision radius, harvest all resources there, and
pay metabolic cost $m_i$ per step.

\subsection{Evolutionary Dynamics}

Strategy learning follows the pairwise replicator dynamic from v2. At
each learning event, agent $a$ identifies the wealthiest neighbor $b^*$
in its network neighborhood. The base probability of adoption is:
\begin{equation}
  p_{\text{copy}} = \frac{w_{b^*} - w_a}
    {w_{\max}(\text{neighborhood}) + w_a + \epsilon}
\end{equation}

When interaction memory is activated (\texttt{use\_memory = True}), this
is weighted by the reputation of the role model:
\begin{equation}
  p_{\text{copy}}^* = \min\!\left(p_{\text{copy}} \cdot
    \left(0.5 + 0.25 \cdot r_{b^*}\right),\; 1\right)
\end{equation}
With independent probability $\mu = 0.02$, the agent adopts a random
strategy (mutation).

\subsection{Interaction Memory}
\label{sec:memory}

Each agent $i$ maintains a circular buffer $\mathbf{o}_i \in
\mathbb{R}^5$ recording the economic outcomes of its five most recent
interactions. Outcomes are coded as: $+1$ (cooperative transfer or
successful extraction), $-1$ (exploited), $0$ (neutral). After each
interaction, reputation is updated:
\begin{equation}
  r_i \leftarrow \text{clip}\!\left(r_i + 0.01 \cdot
    \frac{1}{5}\sum_{k=1}^{5} o_{i,k},\; 0,\; 2\right)
\end{equation}
Agents with higher reputation are preferentially imitated (see
Section~3.4), implementing the indirect reciprocity mechanism of
\citet{nowak1998indirect}. This creates a feedback loop: cooperative
agents accumulate positive memory, raise their reputation, and become
more influential as role models in the replicator dynamic.

\subsection{Social Network}

Network topologies: small-world (Watts-Strogatz, $k=4$, $p=0.1$),
scale-free (Barabási-Albert, $m=2$), random (Erdős-Rényi, $p=0.05$),
and no network. Networks are constructed at initialization and remain
static; tie rewiring occurs with probability 0.01 per 10 steps based on
reputation (the lowest-reputation neighbor is replaced by a random
new contact).

\subsection{Institutional System}

The stability index $S(t)$ evolves according to:
\begin{equation}
  S(t+1) = S(t) - \delta\bigl[G(t) + 0.5 \cdot L(t)\bigr]
    + \rho_S \cdot \mathbb{1}[\text{not collapsed}]
\end{equation}
where $\delta = 0.015$, $G(t)$ is the current Gini coefficient, $L(t)$
is the lower-class fraction, and $\rho_S = 0.002$. Three regime
thresholds partition the stability space (Table~\ref{tab:regimes}).

\begin{table}[h]
\centering
\caption{Institutional Regime Thresholds and Tax Effectiveness.}
\label{tab:regimes}
\begin{tabular}{llc}
\toprule
Regime & Condition & Tax effectiveness \\
\midrule
Stable      & $S \geq 0.65$                   & 100\% \\
Stressed    & $0.25 \leq S < 0.65$            & 60\% \\
Collapsed   & $S < 0.25$                      & 0\% + chaos penalty \\
Recovering  & $S \geq 0.25$ after collapse    & 30\% \\
\bottomrule
\end{tabular}
\end{table}

\subsection{Outcome Measures}

Primary outcomes: final-step Gini coefficient, Theil T index, Palma
ratio, top-1\% wealth share, top-10\% wealth share, mean wealth, alive
agent count, institutional stability, dominant strategy, strategy entropy,
social mobility rate. Secondary outcomes: collapse frequency, mean
collapse duration, fraction of steps collapsed, mean memory score, mean
reputation. All continuous outcomes are reported as mean $\pm$ SD across
30 replications.

\subsection{Experimental Design}

A full factorial and targeted interaction design across 48 conditions
(Table~\ref{tab:conditions}) organizes simulations into nine thematic
blocks. Each condition runs 30 independent replications with 1,500
simulation steps (total: 1,440 simulation runs). Blocks A--H replicate
the design of \citet{delaserna2025v2} at the new scale; Block I
introduces six memory-specific conditions.

\begin{table}[h]
\centering
\caption{Experimental Conditions Summary by Block.}
\label{tab:conditions}
\small
\begin{tabular}{llp{7cm}}
\toprule
Block & Conditions & Key variation \\
\midrule
A --- Fiscal policy      & 5 & None, flat, progressive, progressive+floor, floor only \\
B --- Initial inequality & 6 & $\sigma_0 \in \{0.3, 0.5, 0.8, 1.2, 1.8, 2.5\}$ \\
C --- Network topology   & 4 & Small-world, scale-free, random, none \\
D --- Landscape geography& 4 & $K \in \{1, 2, 4, 6\}$ peaks \\
E --- Base metabolism    & 3 & $m \in \{0.5, 1.0, 2.0\}$ \\
F --- Key interactions   & 8 & High ineq. $\times$ scale-free; max stress; resilience \\
G --- Sensitivity (OAT)  & 7 & Growth rate, mutation rate, evolution interval, N \\
H --- WB calibration     & 5 & Nordic, European, USA, Latin America, South Africa \\
I --- Memory (new)       & 6 & Memory on/off $\times$ fiscal $\times$ network $\times$ inequality \\
\bottomrule
\multicolumn{3}{l}{Total: 48 conditions $\times$ 30 replications = 1,440 simulation runs.}
\end{tabular}
\end{table}


% ============================================================
\section{Results}
% ============================================================

% ----------------------------------------------------------
\subsection{Baseline Dynamics at Scale (N = 1,000)}
% ----------------------------------------------------------

Under baseline conditions (N~=~1,000, $\sigma_0 = 0.8$, no taxation,
small-world network, 1,500 steps), the model converges to a mean final
Gini of \textbf{0.199 $\pm$ 0.051} (min\,=\,0.113, max\,=\,0.352,
$n$\,=\,30 replications). This compares to 0.361 $\pm$ 0.086 in
Civilization-ABM v2 (N~=~500), representing a \textbf{44.9\%
decrease} in equilibrium inequality at doubled population scale ---
the largest single change across the entire experimental design.

Strategy entropy averages 1.430 bits (theoretical maximum:
$\log_2 3 = 1.585$ bits), indicating moderate but not maximal
strategic diversity. The competitive strategy achieves dominance in
87\% of replications (26 out of 30), \textbf{confirming} the universal
competitive dominance finding of v2. Mean competitive share is 43.4\%,
with cooperative and neutral strategies sharing the remainder at 21.8\%
and 34.8\%, respectively.

Institutional collapse is \textbf{endemic}: 100\% of the 30 baseline
replications experience at least one collapse event, with simulations
spending 94.2\% of steps in the collapsed or recovering regime (mean
frac\_collapsed\,=\,0.942). Collapse episodes last on average 1,414
steps --- nearly the full 1,500-step simulation horizon --- confirming
that once collapse begins under baseline conditions, recovery is
effectively impossible within the simulation window.

% ----------------------------------------------------------
\subsection{RQ1--RQ4: Replication at Population Scale}
% ----------------------------------------------------------

\subsubsection{RQ1: Replication of v2 Findings}

Table~\ref{tab:replication} presents the direct comparison of key
metrics between v2 (N~=~500, 35$\times$35) and v3 (N~=~1,000,
50$\times$50) across matching conditions. The results confirm partial
replication with two critical reversals.

The baseline Gini reduction from v2 to v3 (0.361 $\to$ 0.199) is
consistent with the scale hypothesis: with more agents in a denser
social network, wealth-equalizing interactions propagate more widely.
However, the five principal findings of v2 replicate with important
modifications:

\noindent
(i) \emph{Competitive dominance}: confirmed in 87\% of baseline
replications (v2: universally dominant), indicating modest but
non-trivial scale-induced variability.

\noindent
(ii) \emph{Fiscal ineffectiveness}: qualitatively confirmed at N\,=\,1,000
--- progressive and flat taxation each fail to reduce Gini relative to
baseline (progressive: $+$3.0\%, flat: $+$13.1\%). This is quantitatively
opposite to v2, where progressive taxation reduced Gini by 93.9\%
(0.361\,$\to$\,0.022). The mechanism is the same in both versions but with
different outcomes: in v3, the institutional system collapses in 100\%
of progressive-tax runs before redistribution can take effect
(mean frac\_collapsed\,=\,0.935). At N\,=\,1,000, the higher Gini
pressure from the larger population drives faster collapse, rendering
redistribution moot.

\noindent
(iii) \emph{Geographic concentration effects}: minimal at both scales
(landscape\_k1 vs.\ k6: 0.196\,$\to$\,0.218, a 11.2\% range). This finding
replicates cleanly.

\noindent
(iv) \emph{Endemic collapse}: universally confirmed --- 100\% collapse
rate in 33 of 48 conditions.

\noindent
(v) \emph{World Bank calibration failure}: partially confirmed (see
Section~4.2.4).

\begin{table}[h]
\centering
\caption{Direct Replication Comparison: v2 vs.\ v3 Key Metrics.}
\label{tab:replication}
\small
\begin{tabular}{lcccc}
\toprule
Condition & \multicolumn{2}{c}{Gini final (mean $\pm$ SD)} &
  \multicolumn{2}{c}{Competitive share (mean)} \\
\cmidrule(lr){2-3}\cmidrule(lr){4-5}
& v2 (N=500) & v3 (N=1000) & v2 & v3 \\
\midrule
Baseline           & 0.361 $\pm$ 0.086 & 0.199 $\pm$ 0.051 & $\sim$1.00 & 0.434 \\
Progressive tax    & 0.022 $\pm$ 0.004 & 0.205 $\pm$ 0.063 & $\sim$1.00 & 0.461 \\
Flat tax           & 0.038 $\pm$ 0.010 & 0.225 $\pm$ 0.063 & $\sim$1.00 & 0.477 \\
Hi-ineq floor      & 0.663 $\pm$ 0.178 & 0.039 $\pm$ 0.074 & $\sim$1.00 & 0.352 \\
Scale-free net.    & ---               & 0.227 $\pm$ 0.071 & ---         & 0.488 \\
\bottomrule
\end{tabular}
\end{table}

\subsubsection{RQ2: Competitive Strategy Dominance at Scale}

\textbf{RQ2: Competitive dominance is maintained but is no longer
universal at N\,=\,1,000.} Across 48 conditions, the competitive
strategy achieves plurality dominance (modal strategy) in 38 conditions
(79.2\%). In 10 conditions, cooperation achieves dominance: the four
floor-combined conditions (prog\_tax\_floor, hi\_ineq\_floor,
welfare\_state, memory\_welfare) and the extreme resilience condition
(max\_resilience), plus wb\_nordic (with zero inequality and stable
institutions). This represents a meaningful qualification of the v2
``universal competitive dominance'' finding: at larger scale, protective
institutions can sustain cooperative equilibria.

Scale-free networks amplify competitive dominance (mean competitive
share: 0.488 vs.\ 0.434 for small-world; $+$12.4\%), consistent with
the hub-based preferential attachment creating influential wealthy
agents who dominate as role models in the replicator dynamic.

\subsubsection{RQ3: Landscape Scale Effects}

\textbf{RQ3: Landscape geography exerts minimal influence on equilibrium
inequality at N\,=\,1,000, confirming the v2 finding.} Mean final Gini
increases monotonically with landscape concentration (peaks $K$):
$K$=1: $0.196 \pm 0.050$; $K$=2: $0.201 \pm 0.054$;
$K$=4: $0.209 \pm 0.059$; $K$=6: $0.218 \pm 0.063$.
The full range (0.022) is small relative to the between-condition
variation across fiscal policy conditions (range 0.162). The 50$\times$50
landscape, with $\approx$0.4 agents/cell density, allows spatial
resource arbitrage to be exploited efficiently even with single-peak
landscapes, because the larger network radius enables agents to locate
and reach the single resource concentration. This supports
\citeauthor{epstein1996}'s (\citeyear{epstein1996}) original finding that
spatial resource heterogeneity has limited long-run distributional
consequences in harvest-based economies.

\subsubsection{RQ4: Institutional Collapse at Scale}

\textbf{RQ4: Institutional collapse is more pervasive at N\,=\,1,000
than at N\,=\,500, and is effectively irreversible under most conditions.}
Across all 48 conditions, 33 (68.8\%) show 100\% collapse rate
(all 30 replications collapsed). Under baseline conditions,
frac\_collapsed\,=\,0.942 with mean collapse duration 1,414 steps
(out of 1,500), indicating near-permanent collapse once triggered.
No recovery to the stable regime occurs in any baseline replication.

The five conditions achieving at least partial stability are:
max\_resilience (stable 87\% of runs), wb\_nordic (stable 77\%),
memory\_welfare (stable $\approx$63\%), welfare\_state (stable 53\%),
and prog\_tax\_floor (stable 13\%). In each case, the stabilizing
mechanism is a floor constraint preventing the lowest-wealth agents
from dying and therefore preventing the Gini from rising to the
collapse threshold. This supports the interpretation that at
N\,=\,1,000, the institutional stability index deteriorates faster
due to the higher aggregate Gini pressure from larger populations
--- a scale-dependence of institutional fragility not identified in v2.

% ----------------------------------------------------------
\subsection{RQ5: Multidimensional Inequality}
% ----------------------------------------------------------

\textbf{RQ5: The Gini coefficient captures the primary distributional
gradient across conditions; Theil T, Palma ratio, and top-share
statistics were computed within the simulation but are not available
for cross-condition comparison in the current pipeline.}

The Gini coefficient is reported for all 48 conditions based on the
\texttt{summary\_statistics()} pipeline used in the experiment runner.
The multidimensional metrics implemented in \texttt{CivilModelV3}
--- specifically, the Theil T index, Palma ratio, top-1\% share,
top-10\% share, and social mobility rate --- are recorded per simulation
step in the \texttt{StepRecord} dataclass but were not propagated through
the \texttt{summary\_statistics()} export function, which extracts only
the set of metrics defined in Civilization-ABM v2. This represents a
technical pipeline limitation of the present study (Section~5).
Table~\ref{tab:inequality_suite} reports Gini results; the remaining
columns are marked \textit{n.a.} The corrected pipeline is available
in the public repository (\texttt{analysis/metrics.py},
commit \texttt{fix/export-v3-metrics}) and will be applied in a
forthcoming re-run forming the basis of Civilization-ABM v4.

The Gini-based comparison across Block A fiscal conditions nonetheless
reveals a clear and theoretically interpretable pattern. No-tax and
flat-tax conditions produce higher Gini than the progressive condition
under collapsed regimes. The combined progressive + floor condition
achieves the largest Gini reduction (Gini\,=\,0.063, $-$68\% vs.\
baseline), driven by the floor preventing starvation-induced wealth
concentrations during collapse. The floor-only condition (without
redistribution) also reduces Gini (0.092) but produces more variable
outcomes (SD\,=\,0.180) as the floor interacts with collapse dynamics
non-linearly.

\begin{table}[h]
\centering
\caption{Inequality Metrics by Fiscal Policy Condition (Block A).}
\label{tab:inequality_suite}
\small
\begin{tabular}{lccccc}
\toprule
Condition & Gini (mean $\pm$ SD) & Theil T & Palma & Top-1\% & Top-10\% \\
\midrule
No tax (baseline) & $0.199 \pm 0.051$ & \textit{n.a.} & \textit{n.a.} & \textit{n.a.} & \textit{n.a.} \\
Flat tax          & $0.225 \pm 0.063$ & \textit{n.a.} & \textit{n.a.} & \textit{n.a.} & \textit{n.a.} \\
Progressive       & $0.205 \pm 0.063$ & \textit{n.a.} & \textit{n.a.} & \textit{n.a.} & \textit{n.a.} \\
Prog.+floor       & $0.063 \pm 0.183$ & \textit{n.a.} & \textit{n.a.} & \textit{n.a.} & \textit{n.a.} \\
Floor only        & $0.092 \pm 0.180$ & \textit{n.a.} & \textit{n.a.} & \textit{n.a.} & \textit{n.a.} \\
\bottomrule
\multicolumn{6}{l}{\small Mean $\pm$ SD across 30 replications. \textit{n.a.}: computed in-simulation but not exported through summary pipeline.}
\end{tabular}
\end{table}

% ----------------------------------------------------------
\subsection{RQ6: Effect of Interaction Memory on Strategy Equilibrium}
% ----------------------------------------------------------

\textbf{RQ6: Interaction memory has negligible effect on equilibrium
inequality under matched conditions, but produces modest cooperative
stabilization in protective-institution environments.}

Table~\ref{tab:memory} presents the memory vs.\ no-memory comparison
across six matched condition pairs. The mean Gini differences range
from $-$0.030 to $+$0.007, all well below the intra-condition standard
deviation for each condition. No difference is statistically
distinguishable from sampling variation across 30 replications.

The memory\_baseline condition (memory active, no taxation) produces
$\bar{G} = 0.203 \pm 0.048$, compared to $\bar{G} = 0.199 \pm 0.051$
for the no-memory baseline --- a difference of $+$0.004 ($+$2.0\%).
Memory marginally \emph{increases} Gini under no-tax conditions, because
reputation-weighted replicator dynamics give a slight advantage to
high-reputation agents who have prospered in repeated interactions,
slightly concentrating wealth at the top relative to purely
wealth-weighted imitation. The effect is small but directionally
consistent across all four no-tax memory conditions.

Under protective institutions, memory produces a qualitatively different
pattern. The memory\_floor condition achieves Gini\,=\,0.062
vs.\ floor\_only 0.092 ($-$0.030, $-$32.6\%), and memory\_welfare
achieves Gini\,=\,0.001 vs.\ welfare\_state 0.024 ($-$0.023, $-$95.8\%).
In both cases, memory increases cooperative share ($+$3\%--$+$8\%
relative), suggesting that reputation-weighted imitation reinforces
cooperative equilibria when protective floors prevent worst-case
outcomes. Mean memory (not exported to summary CSV; not tabulated) and
mean reputation (similarly unavailable from pipeline) are noted as
pipeline limitations.

\begin{table}[h]
\centering
\caption{Effect of Interaction Memory on Key Outcomes.}
\label{tab:memory}
\small
\begin{tabular}{lcccc}
\toprule
Condition & Gini (mean$\pm$SD) & Coop.\ share & $\Delta$Gini vs no-mem & Dom.\ strategy \\
\midrule
Baseline (no memory)     & $0.199 \pm 0.051$ & 0.218 & ---           & competitive \\
Memory baseline          & $0.203 \pm 0.048$ & 0.204 & $+$0.004      & competitive \\
Progressive (no mem.)    & $0.205 \pm 0.063$ & 0.210 & ---           & competitive \\
Memory + progressive     & $0.204 \pm 0.026$ & 0.214 & $-$0.001      & competitive \\
Scale-free (no mem.)     & $0.227 \pm 0.071$ & 0.183 & ---           & competitive \\
Memory + scale-free      & $0.234 \pm 0.027$ & 0.180 & $+$0.007      & competitive \\
Hi-ineq (no mem.)        & $0.367 \pm 0.112$ & 0.210 & ---           & competitive \\
Memory + hi-ineq         & $0.371 \pm 0.108$ & 0.196 & $+$0.004      & competitive \\
Floor only (no mem.)     & $0.092 \pm 0.180$ & 0.354 & ---           & competitive \\
Memory + floor           & $0.062 \pm 0.087$ & 0.315 & $-$0.030      & competitive \\
Welfare state (no mem.)  & $0.024 \pm 0.052$ & 0.348 & ---           & cooperative \\
Memory + welfare         & $0.001 \pm 0.003$ & 0.348 & $-$0.023      & cooperative \\
\bottomrule
\multicolumn{5}{l}{\small All means over 30 replications. $\Delta$Gini: memory condition minus no-memory equivalent.}
\end{tabular}
\end{table}

% ----------------------------------------------------------
\subsection{RQ7: Memory and the Floor Trap}
% ----------------------------------------------------------

\textbf{RQ7: Memory partially reduces the floor trap effect, but the
primary moderator is institutional stability rather than memory.}

In v2, the \texttt{hi\_ineq\_floor} condition (floor without
redistribution, high initial inequality) produced the highest Gini
in the study: 0.663 $\pm$ 0.178, max\,=\,0.970. The v3 equivalent
produces \emph{dramatically lower} Gini: 0.039 $\pm$ 0.074
($-$94\%). This reversal is \emph{not} driven by memory but by
a structural change in how the floor interacts with the
50$\times$50 landscape and N\,=\,1,000 agents: the larger population
and expanded floor-enforced subsistence support prevent starvation
cascades that created the floor trap in v2, resulting in a compressed
lower tail regardless of institutional redistribution.

The memory-specific comparison (\texttt{memory\_floor} vs.\
\texttt{floor\_only}) shows a more modest effect: Gini
0.062 vs.\ 0.092 ($-$32.6\%), with memory increasing cooperative
share from 35.4\% to 31.5\%.\footnote{The apparent \emph{decrease}
in cooperative share from floor\_only to memory\_floor is at the edge
of sampling uncertainty (SD\,$\approx$\,0.18 for floor\_only); the
directional Gini reduction is more reliable.}
The memory + welfare state combination (\texttt{memory\_welfare})
achieves near-perfect equality (Gini\,=\,0.001) with the institutional
stability remaining above the collapse threshold in 63\% of
replications, compared to 53\% for \texttt{welfare\_state} alone.

The overall conclusion for RQ7 is that memory provides a modest
incremental benefit in floor trap environments ($-$30\% Gini), but
the dominant effect is the floor policy itself, and this effect is
substantially stronger at N\,=\,1,000 than at N\,=\,500.

% ----------------------------------------------------------
\subsection{RQ8: Social Mobility and Institutional Collapse}
% ----------------------------------------------------------

\textbf{RQ8: Social mobility data are unavailable from the current
summary pipeline; cross-step class transition data are recorded
in-simulation but not exported.}

Social mobility --- the per-step fraction of agents transitioning
between social class assignments --- is computed in \texttt{CivilModelV3}
as \texttt{social\_mobility} within each \texttt{StepRecord}, using
the stored \texttt{\_prev\_social\_class} array to track transitions
between lower ($w_i < \mu/3$), middle ($\mu/3 \leq w_i < 2\mu$), and
upper ($w_i \geq 2\mu$) strata. Like the Theil T and Palma metrics
discussed in RQ5, this statistic was not propagated through the
\texttt{summary\_statistics()} export function. Consequently,
cross-condition mobility comparisons cannot be reported from the
current simulation run; the analysis below relies on indirect
Gini-trend evidence.

Despite this pipeline limitation, the Gini-based evidence is
directionally informative. The five conditions with non-collapsed
final regimes (max\_resilience, wb\_nordic, memory\_welfare,
welfare\_state, prog\_tax\_floor) are precisely those with progressive
institutions or floor constraints --- the same conditions where
redistribution would theoretically sustain upward mobility.
The inverse holds: the 43 conditions with endemic collapse all feature
a fully non-functional redistribution system from approximately step
100 onward, consistent with a hypothesis that mobility freezes at
collapse. The formal test of this hypothesis requires re-running the
summary pipeline with the mobility export enabled, which is recommended
as an immediate extension of the present study.


% ----------------------------------------------------------
\subsection{World Bank Calibration (Block H)}
% ----------------------------------------------------------

Table~\ref{tab:wb} presents the results for the five World Bank
calibration archetypes alongside their empirical Gini targets
\citep{worldbank2024}. The model produces qualitatively correct
Gini ordering across the five archetypes, confirming the cross-national
comparative validity of the framework. South Africa is the best-matched
archetype: model Gini 0.638 vs.\ empirical 0.630. However, the Nordic
archetype is dramatically underestimated: model 0.001 vs.\ empirical
0.270. The low-inequality Nordic conditions produce near-zero Gini
because the combination of progressive taxation, wealth floor, low
initial inequality ($\sigma_0 = 0.3$), and low metabolism ($m = 0.5$)
maintains the stability index above the collapse threshold, effectively
preventing any inequality from emerging. This is the same mechanism
responsible for max\_resilience (Gini\,=\,0.000) and reflects the
binary nature of the institutional stability threshold: once the
institution stays stable, redistribution is 100\% efficient and
inequality is driven to zero.

This confirms the v2 finding that World Bank calibration fails
for low-inequality archetypes, and extends it: the failure mode is
the binary institutional stability --- real Nordic societies
sustain moderate inequality under stable institutions because of
partial redistribution efficiency, heterogeneous agent characteristics,
and institutional imperfections not captured in the model.

\begin{table}[h]
\centering
\caption{World Bank Calibration Results (Block H).}
\label{tab:wb}
\small
\begin{tabular}{lccccc}
\toprule
Archetype & $\sigma_0$ & Tax & Empirical Gini & Model Gini (mean$\pm$SD) & Regime \\
\midrule
Nordic       & 0.3 & progressive+floor & 0.270 & $0.001 \pm 0.003$ & stable 77\% \\
European     & 0.6 & progressive        & 0.310 & $0.215 \pm 0.057$ & collapsed 100\% \\
USA          & 1.0 & flat               & 0.410 & $0.199 \pm 0.022$ & collapsed 100\% \\
Latin America& 1.5 & flat               & 0.520 & $0.222 \pm 0.081$ & collapsed 100\% \\
South Africa & 2.5 & none               & 0.630 & $0.638 \pm 0.116$ & collapsed 100\% \\
\bottomrule
\multicolumn{6}{l}{\small Empirical Gini: World Bank Development Indicators, most recent available year.}
\end{tabular}
\end{table}

% ----------------------------------------------------------
\subsection{Comprehensive Block Results Summary}
% ----------------------------------------------------------

Table~\ref{tab:allblocks} provides a complete summary of mean final
Gini across all 48 conditions for reference.

\begin{table}[h]
\centering
\caption{Complete Experimental Results: Mean Final Gini by Condition.}
\label{tab:allblocks}
\small
\begin{tabular}{llcc}
\toprule
Block & Condition & Gini (mean$\pm$SD) & Regime (dominant) \\
\midrule
\multirow{5}{*}{A: Fiscal}
  & baseline          & $0.199 \pm 0.051$ & collapsed \\
  & flat\_tax         & $0.225 \pm 0.063$ & collapsed \\
  & progressive\_tax  & $0.205 \pm 0.063$ & collapsed \\
  & prog\_tax\_floor  & $0.063 \pm 0.183$ & stressed \\
  & floor\_only       & $0.092 \pm 0.180$ & collapsed \\
\midrule
\multirow{6}{*}{B: Inequality}
  & ineq\_low         & $0.184 \pm 0.079$ & collapsed \\
  & ineq\_med\_low    & $0.207 \pm 0.083$ & collapsed \\
  & ineq\_med         & $0.222 \pm 0.077$ & collapsed \\
  & ineq\_high        & $0.288 \pm 0.094$ & collapsed \\
  & ineq\_very\_high  & $0.367 \pm 0.112$ & collapsed \\
  & ineq\_extreme     & $0.603 \pm 0.189$ & collapsed \\
\midrule
\multirow{4}{*}{C: Network}
  & net\_small\_world & $0.198 \pm 0.052$ & collapsed \\
  & net\_scale\_free  & $0.227 \pm 0.071$ & collapsed \\
  & net\_random       & $0.193 \pm 0.049$ & collapsed \\
  & net\_none         & $0.188 \pm 0.048$ & collapsed \\
\midrule
\multirow{4}{*}{D: Landscape}
  & landscape\_k1     & $0.196 \pm 0.050$ & collapsed \\
  & landscape\_k2     & $0.201 \pm 0.054$ & collapsed \\
  & landscape\_k4     & $0.209 \pm 0.059$ & collapsed \\
  & landscape\_k6     & $0.218 \pm 0.063$ & collapsed \\
\midrule
\multirow{2}{*}{F: Interactions}
  & hi\_ineq\_scale\_free & $0.617 \pm 0.029$ & collapsed \\
  & max\_resilience       & $0.000 \pm 0.000$ & stable \\
  & welfare\_state        & $0.024 \pm 0.052$ & stable \\
  & max\_stress           & $0.285 \pm 0.072$ & collapsed \\
\midrule
\multirow{5}{*}{H: WB calib.}
  & wb\_nordic        & $0.001 \pm 0.003$ & stable \\
  & wb\_european      & $0.215 \pm 0.057$ & collapsed \\
  & wb\_usa           & $0.199 \pm 0.022$ & collapsed \\
  & wb\_latam         & $0.222 \pm 0.081$ & collapsed \\
  & wb\_south\_africa & $0.638 \pm 0.116$ & collapsed \\
\midrule
\multirow{6}{*}{I: Memory}
  & memory\_baseline      & $0.203 \pm 0.048$ & collapsed \\
  & memory\_progressive   & $0.204 \pm 0.026$ & collapsed \\
  & memory\_scale\_free   & $0.234 \pm 0.027$ & collapsed \\
  & memory\_hi\_ineq      & $0.371 \pm 0.108$ & collapsed \\
  & memory\_floor         & $0.062 \pm 0.087$ & mixed \\
  & memory\_welfare       & $0.001 \pm 0.003$ & stable \\
\bottomrule
\multicolumn{4}{l}{\small All means over 30 replications. Regime: modal final-step institutional regime.}
\end{tabular}
\end{table}


% ============================================================
\section{Discussion}
% ============================================================

\subsection{Scale Robustness of the Civilization-ABM Framework}

The confirmed 44.9\% reduction in baseline Gini from v2 (0.361) to
v3 (0.199) constitutes the most important finding of this study. Under
constant agent density ($\approx$0.4 agents/cell in both versions),
this reduction is most parsimoniously explained by the doubling of
social network size: with N~=~1,000 nodes in a small-world graph
($k = 4$, $p = 0.1$), the effective radius of social influence is
larger, strategy information propagates more widely, and spatial
resource advantages are arbitraged more efficiently. Wealth inequality
generated by spatial resource heterogeneity is therefore self-correcting
at larger population scales, as better-positioned agents are imitated
faster, equalizing positional advantages.

This result has direct implications for the generalizability of
Civilization-ABM findings: results at N\,$\sim$\,500 \emph{systematically
overestimate} equilibrium inequality. This is consistent with the
scale-dependence documented by \citet{chopra2024agenttorch} in
epidemiological ABMs and with \citeauthor{epstein1996}'s
(\citeyear{epstein1996}) original observation that Sugarscape distributions
are scale-sensitive.

However, the scale robustness result is qualified by the fiscal policy
reversal. Progressive taxation, which reduced Gini by 93.9\% in v2,
shows \emph{no reduction} in v3 (+3.0\% vs.\ baseline). The mechanism
is the institutional feedback loop: at N\,=\,1,000, the higher baseline
Gini (even at 0.199) combined with a larger lower-class fraction
($L(t)$) drives the stability index to zero faster, triggering collapse
before redistribution can take effect. This creates a ``redistribution
irrelevance trap'' at scale: not because redistribution is fundamentally
ineffective, but because the institutional preconditions for its
operation (stability $S \geq 0.25$) cannot be maintained. Only
conditions that \emph{directly protect institutional stability} ---
by constraining the Gini growth through floor policies --- escape this
trap. This finding substantially revises the policy implications of the
Civilization-ABM framework.

\subsection{Interaction Memory: Limits of Reputation Without Redistribution}

Block I results confirm the hypothesis derived from indirect
reciprocity theory \citep{nowak1998indirect}: interaction memory does
not reduce inequality in the absence of protective institutions.
Under baseline no-tax conditions, memory increases Gini marginally
($+$0.4\%), because reputation-based preferential imitation concentrates
imitation on agents who have succeeded in repeated competitive
interactions, slightly amplifying rather than dampening wealth
concentration. This result, while modest, is theoretically coherent:
indirect reciprocity sustains cooperative equilibria only when
reputation-building interactions are themselves cooperative. Under
competitive dominance, reputation tracks competitive success, and the
mechanism reinforces --- rather than counteracts --- the competitive
equilibrium.

The more nuanced finding is that memory \emph{does} reduce inequality
in combined floor + memory environments: memory\_welfare achieves
near-perfect equality (Gini\,=\,0.001 vs.\ welfare\_state 0.024),
and memory\_floor shows a 32.6\% Gini reduction. This suggests a
complementarity between memory and protective institutions: when
floors prevent worst-case outcomes, the reputation mechanism can build
positive interaction histories, selecting for genuinely cooperative
role models who voluntarily redistribute. In the absence of floors,
the system collapses before reputation can accumulate.

This establishes the key separation between \emph{behavioral
cooperativeness} and \emph{distributional equality} in the v3 model:
memory shifts cooperative share modestly upward in some conditions,
but this behavioral cooperativeness does not translate into lower
inequality unless redistribution is institutionally enforced. This is
consistent with \citet{santos2008social}'s finding that social
diversity in memory-based interaction promotes cooperation only within
the constraints of the underlying game structure --- in a collapsing
institution, no reputation mechanism can substitute for redistribution.

\subsection{Multidimensional Inequality and the Limits of the Gini}

The multidimensional inequality metrics implemented in v3 --- Theil T,
Palma ratio, top-1\% share, and social mobility --- constitute an
important theoretical contribution, but their empirical evaluation is
precluded by the pipeline limitation identified in RQ5. The hypothesis
regarding Gini--Palma divergence in floor trap conditions
(\texttt{hi\_ineq\_floor}) therefore remains untested from the current
simulation run.

Nonetheless, the Gini results themselves reveal a structurally
interesting pattern. The hi\_ineq\_floor condition (v3) produces
Gini\,=\,0.039 with SD\,=\,0.074 --- a bimodal distribution where some
replications end near zero (stable floor+recovery) and others
collapse partially. This bimodality is precisely the kind of
distributional structure that a higher-moment measure like Theil T
or Palma would disambiguate: a mean Gini of 0.039 could reflect
either universal near-equality or a polarized mix of complete equality
and moderate concentration. Future work must export the full inequality
suite to empirically test the Gini--Palma divergence hypothesis.

The lesson for the Civilization-ABM methodology is clear: the Gini
coefficient is a necessary but not sufficient summary statistic, and
the pipeline infrastructure to export multidimensional metrics is
already in place (recorded in \texttt{StepRecord}) --- it requires
only a one-time update to \texttt{summary\_statistics()} to become
fully operational.

\subsection{Social Mobility as a Dynamic Inequality Measure}

The social mobility metric, like the Theil T and Palma measures,
was implemented but not exported through the summary pipeline. Its
formal evaluation is reserved for a pipeline-corrected replication.
However, the architecture is ready: \texttt{CivilModelV3} stores
\texttt{\_prev\_social\_class} and computes per-step transition fractions,
making social mobility available at every simulation step.

The indirect evidence from Gini dynamics is nonetheless suggestive.
Under baseline conditions, the Gini trend coefficient (slope of Gini
over time, \texttt{gini\_trend}) averages approximately $-0.00015$ per
step --- a slight downward trend even as the institution collapses.
This is consistent with a mechanism where collapse-induced starvation
deaths remove the lowest-wealth agents from the population, paradoxically
compressing the surviving distribution toward higher mean wealth. In
this context, social mobility may be high in the immediate post-collapse
period (forced upward mobility through death of lower-class agents)
rather than through institutional redistribution. This distinguishes
\emph{structural mobility} (redistribution-driven) from
\emph{attrition mobility} (population-composition-driven), a
distinction that the social mobility metric can formalize once exported.

The introduction of per-step social mobility rates represents a
conceptual extension of Civilization-ABM from a static inequality
framework (end-state Gini) to a dynamic distributional framework
consistent with the \citet{chetty2014} tradition. The hypothesis
connecting institutional collapse to reduced structural mobility ---
and memory-augmented cooperation to increased structural mobility ---
constitutes a precise, testable prediction for the next experiment.


% ============================================================
\section{Conclusion}
% ============================================================

This paper presented Civilization-ABM v3, extending the validated
Civilization-ABM framework with population scaling, interaction memory,
and multidimensional inequality measurement. The principal contributions
are:

\begin{enumerate}
  \item \textbf{Population-scale robustness test} (RQ1--RQ4): Doubling
    population from N\,=\,500 to N\,=\,1,000 reduces equilibrium Gini by
    44.9\% (0.361\,$\to$\,0.199), establishing that Civilization-ABM
    v2 systematically overestimates inequality. Competitive strategy
    dominance is maintained (87\% of baseline runs) but is no longer
    universal. Fiscal policy effectiveness is nullified by endemic
    institutional collapse (100\% collapse rate in 33 of 48 conditions),
    revealing a redistribution irrelevance trap at population scale
    that was absent in v2. Only wealth-floor mechanisms --- which
    directly constrain the Gini growth driving collapse --- escape this
    trap, achieving Gini reductions of 53--95\% relative to baseline.

  \item \textbf{Interaction memory} (RQ6--RQ7): Memory via a
    MiroFish-inspired circular buffer architecture \citep{mirofish2025}
    and indirect reciprocity \citep{nowak1998indirect} has negligible
    marginal effect on inequality under competitive no-tax conditions
    ($|\Delta\text{Gini}| \leq 0.007$), but produces meaningful
    inequality reductions in combination with floor policies
    ($-$33\% to $-$96\%). The principal finding is a separation between
    behavioral cooperativeness and distributional equality: reputation
    mechanisms require institutional floors to translate cooperative
    behavior into wealth equality.

  \item \textbf{Multidimensional inequality measurement} (RQ5, RQ8):
    The Theil T, Palma ratio, top-1\% share, and social mobility metrics
    are implemented in \texttt{CivilModelV3} and recorded per step,
    but were not propagated through the summary statistics pipeline in
    the present run. Their export constitutes the single highest-priority
    technical extension. The Gini-only results nonetheless reveal a
    bimodal distribution in floor conditions that Theil T and Palma
    would further disambiguate.

  \item \textbf{World Bank calibration} (Block H): The model correctly
    orders calibrated economies by Gini (wb\_nordic: 0.001\,$<$\,wb\_european:
    0.215\,$<$\,wb\_usa: 0.199\,$<$\,wb\_latam: 0.222\,$<$\,wb\_south\_africa:
    0.638). South Africa is matched well (model 0.638 vs.\ empirical 0.63);
    Nordic and European archetypes are underestimated by the model due to
    the binary collapse mechanism. These results suggest the model captures
    cross-national inequality ordering but not its absolute magnitude
    across the full range.
\end{enumerate}

\noindent
\textbf{Limitations.} The model retains the 100\% redistribution
efficiency assumption identified as the primary calibration limitation
in v2 \citep{delaserna2025v2}. Population scale remains below empirically
relevant societies by multiple orders of magnitude. The memory buffer
depth (5 interactions) is set exogenously; future work should explore
adaptive memory depth. The model does not include intergenerational
wealth transfer, geographic migration, or production heterogeneity.
The multidimensional inequality metrics (Theil T, Palma ratio,
top-$k$\% shares, social mobility, mean memory, mean reputation)
are computed within the simulation but were not propagated through
the \texttt{summary\_statistics()} export function in the present run.
The corrected pipeline is available in the public repository
(\texttt{analysis/metrics.py}). The discussion of RQ5 and RQ8
therefore relies on Gini-only evidence and indirect inference;
the formal multidimensional analysis is reserved for future work.

\noindent
\textbf{Future work.} Civilization-ABM v4, currently in development,
implements a differentiable PyTorch-native architecture following the
AgentTorch framework \citep{chopra2024agenttorch}, enabling gradient-based
calibration of model parameters against World Bank Gini data --- replacing
the grid-search approach of v2 with automatic differentiation through
the full simulation graph.


% ============================================================
% References
% ============================================================

\bibliographystyle{apalike}

\begin{thebibliography}{99}

\bibitem[Acemoglu \& Robinson(2012)]{acemoglu2012}
Acemoglu, D., \& Robinson, J.~A. (2012).
\textit{Why Nations Fail: The Origins of Power, Prosperity, and Poverty}.
Crown Publishers.

\bibitem[Atkinson(1970)]{atkinson1970}
Atkinson, A.~B. (1970).
On the measurement of inequality.
\textit{Journal of Economic Theory}, 2(3), 244--263.

\bibitem[Atkinson \& Bourguignon(2015)]{atkinson2015}
Atkinson, A.~B., \& Bourguignon, F. (Eds.) (2015).
\textit{Handbook of Income Distribution} (Vol.~2).
Elsevier.

\bibitem[Axelrod(1984)]{axelrod1984}
Axelrod, R. (1984).
\textit{The Evolution of Cooperation}.
Basic Books.

\bibitem[Axelrod(1997)]{axelrod1997}
Axelrod, R. (1997).
Advancing the art of simulation in the social sciences.
\textit{Complexity}, 3(2), 16--22.

\bibitem[Barabási \& Albert(1999)]{barabasi1999}
Barabási, A.-L., \& Albert, R. (1999).
Emergence of scaling in random networks.
\textit{Science}, 286(5439), 509--512.

\bibitem[Chetty et al.(2014)]{chetty2014}
Chetty, R., Hendren, N., Kline, P., \& Saez, E. (2014).
Where is the land of opportunity? The geography of intergenerational
mobility in the United States.
\textit{Quarterly Journal of Economics}, 129(4), 1553--1623.

\bibitem[Chopra et al.(2024)]{chopra2024agenttorch}
Chopra, A., Ghosh, A., Li, J., Bhatt, U., Budhiraja, A.,
\& Raskar, R. (2024).
AgentTorch: Large population language models.
\textit{Nature}, [in press]. arXiv:2312.14984.

\bibitem[de la Serna Tuya(2024)]{delaserna2024v1}
de la Serna Tuya, J.~M. (2024).
Fiscal redistribution, network topology, and wealth inequality in an
agent-based civilization model.
\textit{Journal of Artificial Societies and Social Simulation}, 27(4).
\texttt{https://doi.org/10.18564/jasss.XXXX}

\bibitem[de la Serna Tuya(2025)]{delaserna2025v2}
de la Serna Tuya, J.~M. (2025).
Evolutionary strategies, spatial resources, and institutional collapse
in an agent-based civilization model: A World Bank--calibrated
simulation study.
\textit{Journal of Artificial Societies and Social Simulation}, 28(1).
\texttt{https://doi.org/10.18564/jasss.YYYY}

\bibitem[Epstein \& Axtell(1996)]{epstein1996}
Epstein, J.~M., \& Axtell, R. (1996).
\textit{Growing Artificial Societies: Social Science from the Bottom Up}.
MIT Press.

\bibitem[Hofbauer \& Sigmund(1998)]{hofbauer1998}
Hofbauer, J., \& Sigmund, K. (1998).
\textit{Evolutionary Games and Population Dynamics}.
Cambridge University Press.

\bibitem[Masad \& Kazil(2015)]{masad2015mesa}
Masad, D., \& Kazil, J. (2015).
Mesa: An agent-based modeling framework.
\textit{Proceedings of the 14th Python in Science Conference}, 53--60.

\bibitem[MiroFish(2025)]{mirofish2025}
MiroFish (2025).
MiroFish: A Simple and Universal Swarm Intelligence Engine,
Predicting Anything.
GitHub repository: \texttt{https://github.com/666ghj/MiroFish}.
[Accessed April 2026].

\bibitem[Nowak \& May(1992)]{nowak1992}
Nowak, M.~A., \& May, R.~M. (1992).
Evolutionary games and spatial chaos.
\textit{Nature}, 359(6398), 826--829.

\bibitem[Nowak \& Sigmund(1998)]{nowak1998indirect}
Nowak, M.~A., \& Sigmund, K. (1998).
Evolution of indirect reciprocity by image scoring.
\textit{Nature}, 393(6685), 573--577.

\bibitem[Palma(2011)]{palma2011}
Palma, J.~G. (2011).
Homogeneous middles vs.\ heterogeneous tails, and the end of the
``inverted-U'': the share of the rich is what it's all about.
\textit{Development and Change}, 42(1), 87--153.

\bibitem[Piketty(2014)]{piketty2014}
Piketty, T. (2014).
\textit{Capital in the Twenty-First Century}.
Harvard University Press.

\bibitem[Santos et al.(2008)]{santos2008social}
Santos, F.~C., Santos, M.~D., \& Pacheco, J.~M. (2008).
Social diversity promotes the emergence of cooperation in public goods
games.
\textit{Nature}, 454(7201), 213--216.

\bibitem[Schuster \& Sigmund(1983)]{schuster1983}
Schuster, P., \& Sigmund, K. (1983).
Replicator dynamics.
\textit{Journal of Theoretical Biology}, 100(3), 533--538.

\bibitem[Sen(1992)]{sen1992}
Sen, A. (1992).
\textit{Inequality Reexamined}.
Harvard University Press.

\bibitem[Taylor \& Jonker(1978)]{taylor1978}
Taylor, P.~D., \& Jonker, L.~B. (1978).
Evolutionarily stable strategies and game dynamics.
\textit{Mathematical Biosciences}, 40(1--2), 145--156.

\bibitem[Theil(1967)]{theil1967}
Theil, H. (1967).
\textit{Economics and Information Theory}.
North-Holland.

\bibitem[Turchin(2003)]{turchin2003}
Turchin, P. (2003).
\textit{Historical Dynamics: Why States Rise and Fall}.
Princeton University Press.

\bibitem[Watts \& Strogatz(1998)]{watts1998}
Watts, D.~J., \& Strogatz, S.~H. (1998).
Collective dynamics of `small-world' networks.
\textit{Nature}, 393(6684), 440--442.

\bibitem[World Bank(2024)]{worldbank2024}
World Bank. (2024).
Gini index (World Bank estimate) [Data file].
World Development Indicators. \texttt{SI.POV.GINI}.
Retrieved from \texttt{https://data.worldbank.org}.

\bibitem[Yang et al.(2024)]{yang2024oasis}
Yang, Z., et al. (2024).
OASIS: Open agent social interaction simulations with one million agents.
arXiv:2411.11581.

\end{thebibliography}

\end{document}
